% =====================================================================
% Sample introduction for the dropout-GRPO paper.
% Standalone .tex fragment; assumes the same preamble as
% research_draft_revised.tex (amsmath, hyperref, theorem environments,
% \R, \E, \Loss, \KL, \D, \old macros).
% =====================================================================

\section{Introduction}
\label{sec:intro}

Reinforcement learning from verifier rewards has become the dominant
post-training recipe for reasoning language models on math and code
benchmarks~\citep{shao2024deepseekmath,yu2024dapo}.  In its
\emph{Group Relative Policy Optimization} (GRPO) form, the policy is
updated using a group of $K$ rollouts per prompt and a group-mean
baseline, eliminating the need for a learned value network.  The
estimator is
\begin{equation}
\hat g \;=\; \frac{1}{K} \sum_{k=1}^{K}
  A^{(k)} \, \nabla_\theta \log \pi_\theta\!\big(y^{(k)} \,\big|\, x\big),
\qquad
A^{(k)} = r^{(k)} - \mu_r,
\label{eq:intro-grpo}
\end{equation}
where $r^{(k)}$ is the verifier reward on the $k$-th rollout.  The
construction relies on a single, easily overlooked assumption:
\textbf{the $K$ rollouts must differ.}  When all rollouts agree,
$A^{(k)} \equiv 0$ for every $k$, the gradient vanishes, and the step
contains no learning signal.  In standard token-level decoders,
diversity is supplied for free by stochastic sampling from the next-token
distribution, so this assumption is invisible.

\paragraph{The latent-reasoning regime breaks GRPO.}
A recent line of work replaces discrete chain-of-thought tokens with
continuous \emph{latent thoughts} that are fed back as the next input
embedding for several steps before the model commits to a discrete
answer.  \textsc{Coconut}~\citep{hao2024coconut} is the canonical
instance: starting from a prompt embedding $h_0 = \mathrm{embed}_\theta(x)$,
the model unrolls
\begin{equation}
h_t \;=\; f_\theta(h_{t-1}), \qquad t = 1, \dots, T,
\label{eq:intro-coconut}
\end{equation}
and only after $T$ latent steps does it autoregressively decode answer
tokens $y_1, \dots, y_L \mid h_T$.  Conditional on $\theta$ and $x$, the
latent phase is a \emph{deterministic} function of the parameters.  A
group of $K$ rollouts under the same prompt and the same checkpoint
therefore produces $K$ identical latent trajectories, $K$ identical
$h_T$, and --- under greedy or low-temperature decoding --- $K$ identical
answers.  Equation~\eqref{eq:intro-grpo} is identically zero.  GRPO
applied na\"ively to a Coconut-style policy makes no progress.

The natural rescue is to inject randomness somewhere.  Two options are
unsatisfying.  Raising the answer-token temperature only randomizes the
\emph{decoding} phase, not the latent phase, so the gradient still ignores
where most of the parameters live.  Adding Gaussian noise to $h_t$
introduces stochastic dynamics without any principled reading on what is
being optimized, and burns several hyperparameters before it stabilizes.
Both fixes reduce GRPO from a clean policy-gradient method to an
under-justified search.

\paragraph{Our proposal.}
We argue that the most natural source of stochasticity is the one
already inside the network: \textbf{dropout}.  Concretely, we run each
rollout with dropout \emph{on}, draw one Bernoulli mask $\xi^{(k)}$ per
rollout, hold that mask constant across all $T$ latent recurrence steps,
record the mask, and \emph{replay} the same mask during the policy update.
The augmented policy is
\begin{equation}
\Pi_\theta(\xi, y \mid x)
   \;=\; p(\xi) \, \pi_\theta(y \mid x, \xi),
\label{eq:intro-augmented}
\end{equation}
treating $\xi$ as part of the action.  Three properties follow:

\begin{enumerate}[topsep=2pt,itemsep=2pt]
\item \emph{Exogenous noise.} The mask distribution $p(\xi)$ does not
      depend on $\theta$, so $\nabla_\theta \log p(\xi) = 0$ and the
      score function reduces to
      $\nabla_\theta \log \pi_\theta(y \mid x, \xi)$, differentiable
      through the unrolled latent recurrence.
\item \emph{Variational dropout interpretation.}  Holding one mask
      constant across all $T$ latent steps matches the variational
      construction of~\citet{gal2016dropout}: each rollout is a single
      posterior sample $\tilde\theta(\xi) = \theta \odot m(\xi)/q$ from
      a structured weight distribution $q_\theta(\theta')$.  The marginal
      policy $\bar\pi_\theta(y \mid x) = \E_{\xi}[\pi_\theta(y\mid x,\xi)]$
      is then a true Bayesian model average, and the GRPO estimator
      targets its expected reward directly.
\item \emph{Common-random-numbers coupling.}  Replaying the
      rollout-time mask at update time means the only difference between
      the rollout-time and update-time latent trajectories is the
      parameter step $\theta - \theta_\old$ itself.  The PPO importance
      ratio is identically $1$ at $\theta = \theta_\old$ and the latent
      state $h_T$ at update time exactly reproduces the one that was
      graded, eliminating a variance source that a fresh-rollout
      estimator pays.
\end{enumerate}

\paragraph{What we prove.}
We give an inductive argument that the latent Jacobian
$\partial h_T / \partial \theta$ is well defined at every depth $T$
under standard Transformer smoothness, justifying ordinary
backpropagation through the unrolled recurrence
(Lemma~\ref{lem:induction}).  We then show that the dropout-GRPO
surrogate gradient with advantage $A^{(k)} = r^{(k)} - \mu_r$ and
replayed masks satisfies
\begin{equation}
\nabla_\theta \Loss^{\mathrm{surr}}(\theta)\Big|_{\theta = \theta_\old}
  \;=\;
\Big(1 - \tfrac{1}{K}\Big)\,\nabla_\theta J(\theta)\Big|_{\theta_\old},
\label{eq:intro-unbiased}
\end{equation}
where $J(\theta) = \E[r]$ is the expected reward of the marginal policy
$\bar\pi_\theta$ (Theorem~\ref{thm:unbiased},
Corollary~\ref{cor:marginal}).  Crucially, the proof requires that
\emph{no within-group statistic of the rewards beyond the mean enter the
advantage.}  In particular, the standard-deviation normalization
$1/(\sigma_r + \delta)$ that appears in the original GRPO formulation
introduces a coupling between $s^{(k)}$ and every other $r^{(j)}$ via a
nonlinear function of all rewards, which breaks the control-variate
identity and biases the estimator.  This was independently noted
by~\citet{liu2024drgrpo} and is the reason we drop the $\sigma_r$
normalization in this work.

We further establish that mask replay couples the rollout-time and
update-time score-function estimators, yielding a common-random-numbers
variance reduction over a fresh-rollout baseline that scales with the
square of the trust-region radius
(Proposition~\ref{prop:crn}).  Because PPO clipping already constrains
$\|\theta - \theta_\old\|$ to be small, this is the dominant variance
benefit during a typical mini-epoch.

\paragraph{Implementation.}
Realizing the shared-mask construction at the level of an LLM forward
pass is non-trivial.  We capture the CUDA + CPU RNG state at the start of
each rollout and restore it before every latent recurrence step, so that
constant-shape dropout layers (residual / MLP / embedding) draw
bit-identical Bernoulli masks at every step.  Variable-shape attention
dropout, whose mask grows with the kv-cache length, is bit-identical on
the shape-overlap prefix in PyTorch's Philox kernel; a runtime auditor
verifies this.  At update time we regenerate $\xi^{(k)}$ from a stored
seed (a few bytes per rollout), keeping mask storage at
$\mathcal{O}(K \cdot \mathrm{seed})$ rather than
$\mathcal{O}(K \cdot T \cdot D)$.  The result is reproducible to the bit:
the policy log-prob at update time matches its rollout-time value
exactly when $\theta = \theta_\old$, which is the precondition for the
PPO importance ratio to start at $1$.

\paragraph{Practical refinements.}
Dropout-GRPO trains stably on math benchmarks once four implementation
details are in place, each addressing a specific empirical pathology:
\begin{itemize}[topsep=2pt,itemsep=2pt]
\item A \emph{Huberized} $k_3$ KL estimator that keeps tail gradients
      bounded but \emph{nonzero} outside the trust region, replacing the
      common log-ratio clamp whose autograd derivative vanishes beyond
      its bounds.  Tokens whose log-ratio drifts past the clamp boundary
      are then a leading pre-collapse signature; the Huberized form
      retains a restoring force.
\item An \emph{annealed} KL coefficient $\beta(t) = \beta_0\,\eta_t/\eta_{\max}$
      proportional to the learning rate.  With cosine lr decay, holding
      $\beta$ fixed pins the policy at an equilibrium where the
      regularizer pull and the policy gradient cancel; coupling
      $\beta$ to $\eta_t$ shifts that equilibrium outward as training
      progresses.
\item \emph{Sequence-mean} loss aggregation rather than token-mean,
      so that each rollout contributes equally regardless of length.
      Token-mean implicitly weights long (often incorrect) responses
      more, distorting the gradient direction on benchmarks where
      response length correlates with correctness.
\item A \emph{group-level accuracy filter}
      $\alpha_{\min} \le \mu_r \le \alpha_{\max}$ that drops trivially
      easy and trivially hard groups before any backward pass, with
      collective DDP coordination so that empty global batches skip the
      optimizer step entirely.  This concentrates compute on the
      marginal-difficulty band where signal-to-noise is highest, and
      acts as an emergent online curriculum: the band shifts as the
      policy improves.
\end{itemize}

\paragraph{Empirical evidence.}
On \textsc{GSM8K} fine-tuning of a Coconut-style latent-reasoning model,
dropout-GRPO with mask replay achieves \emph{[insert headline metric:
$x.y\%$ pass\@1, vs.\ $x.y\%$ for the deterministic-rollout baseline
which fails to leave its initialization]}.  Ablations isolating the
shared-mask construction, the dropped $\sigma_r$ normalization, and the
four refinements above confirm that each contributes measurably to
final accuracy and to training stability; full results appear in
Section~\ref{sec:experiments}.

\paragraph{Contributions.}
\begin{enumerate}[topsep=2pt,itemsep=2pt]
\item We identify why GRPO fails on latent-reasoning policies and
      propose dropout-as-stochasticity with mask replay as a minimal
      and theoretically clean fix.
\item We prove unbiasedness of the dropout-GRPO surrogate gradient under
      the mean-only advantage $A^{(k)} = r^{(k)} - \mu_r$, give an
      explicit common-random-numbers variance-reduction argument over a
      fresh-rollout baseline, and connect the construction to the
      Bayesian model-average reading of variational dropout.
\item We document four implementation refinements (Huberized KL,
      $\beta$-annealing, sequence-mean aggregation, group accuracy
      filtering with DDP-coordinated skips) that are required in
      practice and clarify which carry theoretical weight versus which
      are empirical stabilizers.
\item We release a reference implementation
      (\texttt{coconut\_variational.py},
      \texttt{run\_grpo\_dropout\_variational.py}) that realizes
      bit-identical mask replay across the entire latent recurrence,
      including a runtime auditor for variable-shape attention dropout.
\end{enumerate}

The remainder of the paper is organized as follows.
Section~\ref{sec:coconut-recurrence} fixes notation for the
dropout-perturbed Coconut recurrence and the augmented policy.
Section~\ref{sec:method} describes the dropout-GRPO procedure.
Section~\ref{sec:revision} contains the well-definedness lemma, the
unbiasedness theorem, the CRN proposition, and the Bayesian
model-average reading.  Section~\ref{sec:impl} details the four
implementation refinements.  Section~\ref{sec:experiments} reports
empirical results.  We close with practical caveats and open
questions.
